
\setcounter{paragraph}{0}
\setcounter{ex}{0} 
\PhanI
\begin{ex} %[3P4Y3-1][Loai1] 
\nguon{ĐH - 2010} Sóng điện từ
\choice
{là sóng dọc hoặc sóng ngang}
{\True là điện từ trường lan truyền trong không gian}
{có thành phần điện trường và thành phần từ trường tại một điểm dao động cùng phương}
{không truyền được trong chân không}
\loigiai{Sóng điện từ là điện từ trường lan truyền trong không gian}
\end{ex} \haidong{20}
\begin{ex} %[3P4B3-1][Loai1] 
Phát biểu nào sau đây \textbf{sai}? Sóng điện từ và sóng cơ
\choice
{đều tuân theo quy luật giao thoa}
{đều mang năng lượng}
{đều tuân theo quy luật phản xạ}
{\True đều truyền được trong chân không}
\loigiai{
Sóng cơ không truyền được trong chân không}
\end{ex} \haidong{20}
\begin{ex} %[3P4B4-1][Loai2] 
Một người đang dùng điện thoại di động có thể thực hiện được cuộc gọi. Lúc này điện thoại phát ra
\choice
{tia tử ngoại}
{bức xạ gamma}
{tia Rơn-ghen}
{\True sóng vô tuyến}
\loigiai{
Điện thoại phát ra sóng vô tuyến.
}
\end{ex} \haidong{20}
\begin{ex} %[3P4Y4-1][Loai3] 
Sóng nào sau đây dùng được trong vô tuyến truyền hình trên mặt đất?
\choice
{Sóng dài}
{Sóng trung}
{\True Sóng ngắn}
{cả A, B, C}
\loigiai{
Sóng ngắn được dùng trong vô truyến truyền hình trên mặt đất.
}
\end{ex} \haidong{20}
\begin{ex} %[3P4Y4-1][Loai3]
Sóng điện từ \textbf{không} bị phản xạ ở tầng điện li là
\choice
{\True sóng cực ngắn}
{sóng ngắn}
{sóng trung}
{sóng dài}
\loigiai{ Sóng điện từ không bị phản xạ ở tầng điện li là sóng cực ngắn.
}
\end{ex} \haidong{20}
\begin{ex} %[3P4B4-1][Loai3]
\nguon{ĐH - 2015} Ở Trường Sa, để có thể xem các chương trình truyền hình phát sóng qua vệ tinh, người ta dùng anten thu sóng trực tiếp từ vệ tinh, qua bộ xử lí tín hiệu rồi đưa đến màn hình. Sóng điện từ mà anten thu trực tiếp từ vệ tinh thuộc loại
\choice
{sóng trung}
{\True sóng cực ngắn}
{sóng dài}
{sóng ngắn}
\loigiai{
Sóng truyền qua vệ tinh là sóng cực ngắn.
}
\end{ex} \haidong{20}
\begin{ex} %[3P4B4-1][Loai3]
Theo thứ tự tăng dần về tần số của các sóng vô tuyến, sắp xếp nào sau đây đúng?
\choice
{Sóng cực ngắn, sóng ngắn, sóng trung, sóng dài}
{Sóng dài, sóng ngắn, sóng trung, sóng cực ngắn}
{Sóng cực ngắn, sóng ngắn, sóng dài, sóng trung}
{\True Sóng dài, sóng trung, sóng ngắn, sóng cực ngắn}
\loigiai{
Bước sóng càng dài thì tần số càng nhỏ.
}
\end{ex} \haidong{20}
\begin{ex} %[3P4Y4-1][Loai5]
Trong chân không, khi càng tăng tần số của nguồn phát sóng điện từ thì
\choice
{sóng điện từ truyền càng nhanh}
{khả năng đâm xuyên của sóng điện từ càng giảm}
{năng lượng sóng điện từ càng giảm}
{\True bước sóng của sóng điện từ càng giảm}
\loigiai{
Tần số điện từ càng tăng, bước sóng càng giảm.
}
\end{ex} \haidong{20}
\begin{ex} %[3P4B4-1][Loai12]
Hoạt động nào sau đây là kết quả của việc truyền thông tin liên lạc bằng sóng vô tuyến?
\choice
{Xem phim từ đầu đĩa $DVD$}
{\True Xem thời sự truyền hình qua vệ tinh}
{Trò chuyện bằng điện thoại bàn}
{Xem phim từ truyền hình cáp}
\loigiai{
Hoạt động xem thời sự truyền hình qua vệ tinh là kết quả của việc truyền thông tin liên lạc bằng sóng vô tuyến.
}
\end{ex} \haidong{20}
























\begin{ex}%[3P5Y8-1][Loai1]
\nguon{QG - 2017} Tính chất nổi bật của tia hồng ngoại là
\choice
{\True có tác dụng nhiệt rất mạnh}
{gây ra hiện tượng quang điện ngoài ở kim loại}
{không bị nước và thủy tinh hấp thụ}
{có khả năng đâm xuyên rất mạnh}
\loigiai{
Tính chất nổi bật của tia hồng ngoại là có tác dụng nhiệt rất mạnh.
}
\end{ex} \haidong{20}
\begin{ex} 
\nguon{QG - 2019} Trong chân không, bức xạ có bước sóng nào sau đây là bức xạ hồng ngoại?
\choice
{500 nm}
{350 nm}
{\True 850 nm}
{700 nm}
\loigiai{
Tia hồng ngoại là bức xạ điện từ có bước sóng dài hơn bước sóng của ánh sáng đỏ (760 nm) và ngắn hơn bước sóng của tia viba (1 mm). 
}
\end{ex} \haidong{20}
\begin{ex} 
\nguon{QG - 2019} Trong chân không, bức xạ có bước sóng nào sau đây là tia tử ngoại?
\choice
{450 nm}
{\True 120 nm}
{750 nm}
{920 nm}
\loigiai{Tia tử ngoại có bước sóng từ khoảng $10\ nm$ đến dưới $380\ nm$.
}
\end{ex} \haidong{20}
\begin{ex}%[3P5Y8-1][Loai1]
Nguồn bức xạ nào sau đây \textbf{không} phát ra tia tử ngoại?
\choice
{Mặt trời}
{Đèn hơi thủy ngân}
{Hồ quang điện}
{\True Ngọn nến}
\loigiai{
%Ngọn nến không phát ra tia tử ngoại.
}
\end{ex} \haidong{20}
\begin{ex}%[3P5Y8-1][Loai1]
\nguon{QG - 2016} Tầng ôzôn là tấm “áo giáp” bảo vệ cho người và sinh vật trên mặt đất khỏi bị tác dụng hủy diệt của
\choice
{\True tia tử ngoại trong ánh sáng Mặt Trời}
{tia đơn sắc màu đỏ trong ánh sáng Mặt Trời}
{tia đơn sắc màu tím trong ánh sáng Mặt Trời}
{tia hồng ngoại trong ánh sáng Mặt Trời}
\loigiai{
Tầng ôzôn ngăn không cho các tia tử ngoại trong ánh sáng Mặt Trời tác động tới người và sinh vật trên mặt đất.
}
\end{ex} \haidong{20}
\begin{ex} %[3P5Y8-1][Loai1]
\nguon{QG - 2020} Phát biểu nào sau đây là \textbf{sai}?
\choice
{\True Tia X có bước sóng lớn hơn bước sóng của ánh sáng đỏ}
{Tia X làm ion hóa không khí}
{Tia X có khả năng đâm xuyên}
{Tia X có bước sóng nhỏ hơn bước sóng của ánh sáng tím}
\loigiai{ Tia X có bước sóng nhỏ hơn bước sóng của ánh sáng đỏ}
\end{ex} \haidong{20}
\begin{ex}%[3P5Y8-1][Loai1]
\nguon{QG - 2016} Tia X \textbf{không} có ứng dụng nào sau đây?
\choice
{Chữa bệnh ung thư}
{Tìm bọt khí bên trong các vật bằng kim loại}
{Chiếu điện, chụp điện}
{\True Sấy khô, sưởi ấm}
\loigiai{
Tia X không có ứng dụng “sấy khô, sưởi ấm”.
}
\end{ex} \haidong{20}
\begin{ex} 
\nguon{QG - 2019} Chiếu điện và chụp điện trong các bệnh viện là ứng dụng của
\choice
{tia $\alpha$}
{tia tử ngoại}
{tia hồng ngoại}
{\True tia X}
\loigiai{
Tia X dùng trong chiếu điện và chụp điện trong bệnh viện.
}
\end{ex} \haidong{20}
\begin{ex}%[3P5Y8-1][Loai1]
Thân thể con người bình thường có thể phát ra được bức xạ nào dưới đây?
\choice
{Tia X}
{Ánh sáng nhìn thấy}
{\True Tia hồng ngoại}
{Tia tử ngoại}
\loigiai{
Cơ thể người có thể phát ra tia hồng ngoại.
}
\end{ex} \haidong{20}















\setcounter{paragraph}{0}
\setcounter{ex}{0} 
\PhanII
\begin{ex}
Sóng điện từ trải rộng từ dải vô tuyến dùng trong thông tin liên lạc tới tia X, tia $\gamma$ ứng dụng trong y học và vật lí hạt nhân.
\choiceTF[t]
{Một sóng vô tuyến có tần số $10^{8} \ Hz$ được truyền trong không trung với tốc độ $3.10^{8} \ m/s$. Bước sóng của sóng đó là $1,5 \ m$}
{Sóng điện từ có bước sóng $7 . 10^{-9} \ m$ thuộc tia hồng ngoại}
{\True Sóng điện từ có bước sóng $2.10^{-8} \ m$ thuộc tia tử ngoại}
{\True Tia X có tần số trong khoảng từ $10^{17} \ Hz$ đến $10^{19} \ Hz$}
\loigiai{
\begin{itemchoice}
\itemch  $\lambda=\dfrac{v}{f}=\dfrac{3.10^{8}}{10^{8}}=3 \ m$
\itemch Sóng điện từ có bước sóng $7.10^{-9} \ m$ thuộc tia X. Vùng tia X có bước sóng từ $10^{-9} \ m$ đến $10^{-12} \ m$
\itemch Sóng điện từ có bước sóng $2 . 10^{-8} \ m$ thuộc tia tử ngoại. Vùng tia tử ngoại có bước sóng từ $0,38 . 10^{-6} \ m$ đến  $10^{-9} \ m$
\itemch 
\end{itemchoice}
}
\end{ex} \haidong{20}
\begin{ex}
Tia hồng ngoại (IR) là vùng bức xạ điện từ mang năng lượng nhiệt, nằm ngay sau ánh sáng đỏ trong phổ  sóng điện từ.
\choiceTF[t]
{\True Tia hồng ngoại được ứng dụng để sưởi ấm}
{Tia hồng ngoại là bức xạ nhìn thấy được, có bước sóng lớn hơn bước sóng của ánh sáng đỏ}
{Các vật phát ra tia hồng ngoại phải có nhiệt độ trên $0^\circ C$. Ở các nhiệt độ cao các vật có thể phát cả tia hồng ngoại và ánh sáng nhìn thấy}
{\True Tia hồng ngoại cũng giống như sóng vô tuyến điện và ánh sáng nhìn thấy, đều có cùng bản chất là các sóng điện từ ở các dải tần số khác nhau}
\loigiai{
\begin{itemchoice}
\itemch
\itemch Tia hồng ngoại là bức xạ không nhìn thấy
\itemch Các vật có nhiệt độ lớn hơn nhiệt độ môi trường thì phát tia hồng ngoại
\itemch
\end{itemchoice}
}
\end{ex} \haidong{20}
\begin{ex}
Tia hồng ngoại là bức xạ điện từ mang năng lượng nhiệt, có bước sóng lớn hơn ánh sáng đỏ và được ứng dụng rộng rãi trong đời sống.
\choiceTF[t]
{\True Tia hồng ngoại là những bức xạ không nhìn thấy, có bước sóng lớn hơn bước sóng của ánh sáng đỏ}
{Chỉ có các vật ở nhiệt độ thấp mới phát ra tia hồng ngoại}
{\True Tác dụng nổi bật nhất của tia hồng ngoại là tác dụng nhiệt}
{\True Tia hồng ngoại có bản chất là sóng điện từ}
\loigiai{
\begin{itemchoice}
\itemch
\itemch Các vật có nhiệt độ lớn hơn nhiệt độ môi trường thì phát tia hồng ngoại
\itemch
\itemch
\end{itemchoice}
}
\end{ex} \haidong{20}
\begin{ex}
Tia tử ngoại (UV) là thành phần bức xạ điện từ có bước sóng ngắn hơn ánh sáng tím, năng lượng lớn và được ứng dụng trong nhiều lĩnh vực từ y sinh đến vật liệu.
\choiceTF[t]
{Tia tử ngoại dùng để kiểm tra hành lý của hành khách đi máy bay}
{\True Tia tử ngoại là bức xạ không nhìn thấy được có bước sóng nhỏ hơn bước sóng của ánh sáng tím}
{\True Đèn thủy ngân là một trong những nguồn phát ra tia tử ngoại}
{\True Người ta tìm ra tia tử ngoại là nhờ khả năng ion hóa chất khí}
\loigiai{
\begin{itemchoice}
\itemch
\itemch Tia X dùng để kiểm tra hành lý của hành khách đi máy bay.
\itemch
\itemch
\itemch
\end{itemchoice}
}
\end{ex} \haidong{20}


























